\begin{explanationbox}
\par\noindent\begin{minipage}{\linewidth}
\textbf{\textcolor{CExplanation}{三块底不同，高却相同}}
\PIncenterFigure{2}{4}{6}{1.05}{a}{b}{c}{1}{蓝块 $IBC$、橙块 $ICA$、绿块 $IAB$：\\[3pt]对应底为 $a,b,c$，蓝色垂线段的长度都为 $r$。}
\end{minipage}\par
只看三块的底与高：每块面积都含有共同因子 $r/2$，所以合并面积时只需相加底边。

\textbf{\textcolor{CExplanation}{把图形翻译成算式}}
\[
\underbrace{S_{\triangle IBC}}_{\frac12 ar}
+\underbrace{S_{\triangle ICA}}_{\frac12 br}
+\underbrace{S_{\triangle IAB}}_{\frac12 cr}
= S.
\]
\[
 S=\frac12r(a+b+c)=rp,\qquad r=\frac Sp.
\]

\medskip
\textbf{\textcolor{CExplanation}{做题时的识别信号与作用}}
\begin{description}[style=nextline,leftmargin=2.8em,labelsep=0.5em,itemsep=0.45em]
  \item[\textbf{信号一：给了面积和周长，求内切圆半径。}]
  直接 $r=2S/L$。\textbf{省去：}作角平分线、求内心坐标、再求点到直线距离。
  \item[\textbf{信号二：给了半径和周长，求面积。}]
  直接 $S=rL/2$。\textbf{省去：}寻找哪条边作底以及对应的高。
  \item[\textbf{信号三：给了三边，求半径。}]
  先判断是不是直角三角形；是就用直角边求面积，
  否则可用已学的面积方法（例如海伦公式），最后用 $r=S/p$。
  \item[\textbf{信号四：出现内心连接顶点、求面积比。}]
  三块小三角形\textbf{高相同}，面积比直接为 $a:b:c$，\textbf{不必先求面积或半径}。
  \item[\textbf{信号五：给了三边和内切圆半径，核对图形。}]
  先算 $S=rp$，再和勾股定理、海伦公式得到的面积比较，检查条件能否同时成立。
\end{description}

\medskip
\textbf{\textcolor{CExplanation}{特殊三角形：先约定摆放方式，再写内心坐标}}

以下均用大写 $X,Y$ 表示坐标轴，避免与切线段长 $x,y,z$ 混淆。
\textbf{同一三角形平移或旋转后，内心坐标会改变，但半径和切线段长不变。}
这里的“等腰直角三角形”指两条直角边相等；等边三角形三个角均为 $60^\circ$，不会同时是直角三角形。

\medskip
\par\noindent\begin{minipage}{\linewidth}
\textbf{\textcolor{CExplanation}{情形一：等边三角形——内心、重心、外心、垂心重合}}

设边长为 $\ell>0$，取
\[
 B=(0,0),\quad C=(\ell,0),\quad
 A=\left(\frac\ell2,\frac{\sqrt3\ell}{2}\right).
\]
\PSpecialTriangle{0}
\end{minipage}\par
由对称性，内心在底边中垂线上，且四心重合。设高为 $h=\sqrt3\ell/2$，则
\[
 \boxed{r=\frac{\sqrt3\ell}{6}=\frac h3},\qquad
 \boxed{I=\left(\frac\ell2,\frac{\sqrt3\ell}{6}\right)}.
\]
\textbf{为什么是高的三分之一？}面积 $S=\ell h/2$，半周长 $p=3\ell/2$，
所以 $r=S/p=h/3$。内心到顶点 $A$ 的距离为 $AI=h-r=2h/3=2r$。
三个切点都是各边中点，因此
\[
 \boxed{x=y=z=\ell/2}.
\]

\medskip
\par\noindent\begin{minipage}{\linewidth}
\textbf{\textcolor{CExplanation}{情形二：直角三角形——坐标等于两个正的垂直距离}}

把直角顶点 $B$ 放在原点，两条直角边沿正坐标轴：
\[
 B=(0,0),\quad C=(u,0),\quad A=(0,v),\qquad u,v>0,
\]
斜边 $w=\sqrt{u^2+v^2}$。图取 $u=3,v=4,w=5$，表示一般直角三角形的结构。
\PSpecialTriangle{1}
\end{minipage}\par
内心在第一象限，到两条坐标轴的距离都为 $r$；
所以横、纵坐标分别都是 $r$，内心落在直角的角平分线 $Y=X$ 上：
\[
 \boxed{I=(r,r)},\qquad
 \boxed{r=\frac{uv}{u+v+w}=\frac{u+v-w}{2}}.
\]
\textbf{半径公式的来由：}$S=uv/2$，故 $r=S/p=uv/(u+v+w)$；
再利用 $w^2=u^2+v^2$，有
\[
 (u+v-w)(u+v+w)=(u+v)^2-w^2=2uv,
\]
便得到后一种写法。\textbf{从直角顶点 $B$ 引出的两条切线段恰好等于半径}，即 $BD=BF=y=r$。
沿用本篇 $x,y,z$ 的顶点对应，另有
\[
 x=\frac{v+w-u}{2},\qquad z=\frac{u+w-v}{2}.
\]
注意：一般直角三角形的内心\textbf{不是斜边中点}，斜边中点是外心。

\medskip
\par\noindent\begin{minipage}{\linewidth}
\textbf{\textcolor{CExplanation}{情形三：等腰直角三角形——半径只需一个边长}}

在上一种摆放方式中令 $u=v=t>0$，即
\[
 B=(0,0),\quad C=(t,0),\quad A=(0,t),\quad AC=\sqrt2t.
\]
\PSpecialTriangle{2}
\end{minipage}\par
内心在对称轴 $Y=X$ 上。把 $u=v=t,w=\sqrt2t$ 代入直角三角形公式：
\[
 \boxed{r=\frac{2-\sqrt2}{2}t},\qquad
 \boxed{I=\left(\frac{2-\sqrt2}{2}t,\frac{2-\sqrt2}{2}t\right)}.
\]
此时 $p=(2+\sqrt2)t/2$，所以各顶点的切线段长为
\[
 \boxed{x=z=\frac{t}{\sqrt2}},\qquad \boxed{y=r=\frac{2-\sqrt2}{2}t}.
\]
\textbf{若给的是斜边 $w$：}先用 $t=w/\sqrt2$，也可直接记
\[
 \boxed{r=\frac{\sqrt2-1}{2}w},\qquad \boxed{x=z=w/2}.
\]
只有两条直角边相等时才能用本小节的单边长公式；普通直角三角形用上一节的 $u,v,w$ 公式。

\medskip
\textbf{\textcolor{CExplanation}{数值回验：看清“边长”指哪条边}}

等边三角形边长为 $6$ 时，$I=(3,\sqrt3)$、$r=\sqrt3$，三组切线段长均为 $3$。
直角三角形按图取 $u=3,v=4$ 时，$I=(1,1)$、$r=1$，$(x,y,z)=(3,1,2)$。
等腰直角三角形若直角边 $t=2$，则 $I=(2-\sqrt2,2-\sqrt2)$、$r=2-\sqrt2$，
而 $x=z=\sqrt2$，不是把 $t$ 当成斜边代入。

\medskip
\textbf{\textcolor{CExplanation}{容易忽略的一点}}

“到三边距离相等”一定要和“\textbf{点在三角形内部}”一起理解。
三角形外部还有到三条边所在直线距离相等的点，因此不能只凭“等距”就说任意位置的点都是内心。
\end{explanationbox}
